\chapter{线性响应理论}
\renewcommand{\currentchapterpath}{chapters/ch05/}

\begin{flushright}
\textit{“见微而知著。” ——典出《吕氏春秋》等典籍所体现的思想
\footnote{弱外场下的线性响应，正是由“微”见“著”：从平衡涨落读出宏观可观测量。}}
\end{flushright}

实验上大量谱学与输运测量处于弱扰动情形。线性响应理论给出：响应函数 $=$ 平衡态推迟对易子关联。本章结合 Giuliani--Vignale 第 3--5 章与 Mahan 第 8 章：Kubo 公式、Lindhard/RPA/局域场，以及电流--电流关联给出的电导率。


\chapterinput{perturbation}
\chapterinput{response}
\chapterinput{periodic}
\chapterinput{spectral}
\chapterinput{density_response}
\chapterinput{lindhard}
\chapterinput{rpa}
\chapterinput{local_field}
\chapterinput{conductivity}
